\par\addvspace{12pt}\Needspace{9\baselineskip}

\section{反常积分*}

% ITEM: calculus-example-084
\solutionitem{calculus-example-084}{例5.84 *}
讨论\(\int_0^{+\infty}xe^{-x}\,dx\)的敛散性并求值。再考察\(f(x)=x\)：对称截取区间的积分极限，能否说明其在全实轴上的反常积分收敛？
\solutionanswer
截断积分为\(\int_0^Rxe^{-x}dx=1-(R+1)e^{-R}\to1\)，故收敛且值为\(1\)。\(f=x\)的对称截断积分恒为零，但正半轴积分趋\(+\infty\)，负半轴积分趋\(-\infty\)。两侧必须分别收敛，才有全实轴反常积分；对称抵消的极限不能代替它。


% ITEM: calculus-exercise-58-01
\solutionitem{calculus-exercise-58-01}{习题5.58（1） *}
判断\(\int_0^1\ln x\,dx\)、\(\int_1^{+\infty}dx/(1+x^2)\)的敛散性，收敛时求值。
\solutionanswer
\(\int_\varepsilon^1\ln xdx=-1-\varepsilon\ln\varepsilon+\varepsilon\to-1\)，所以第一式收敛于\(-1\)。第二式为\(\lim_R[\arctan x]_1^R=\pi/4\)。


% ITEM: calculus-exercise-58-02
\solutionitem{calculus-exercise-58-02}{习题5.58（2） *}
分别讨论\(\int_0^1dx/(x\sqrt{x})\)与\(\int_1^{+\infty}dx/(x\sqrt{x})\)。
\solutionanswer
都是\(x^{-3/2}\)。零点端指数大于\(1\)，发散至\(+\infty\)；无穷端指数大于\(1\)，收敛，值为\(2\)。


% ITEM: calculus-exercise-58-03
\solutionitem{calculus-exercise-58-03}{习题5.58（3） *}
说明为什么不能直接把\(\ln|x|\)在\(-1,1\)两端的值相减，就断定\(\int_{-1}^1dx/x=0\)。
\solutionanswer
被积函数在内部\(0\)无界，须分别取\(\int_{-1}^0dx/x\)、\(\int_0^1dx/x\)，二者均发散。\(\ln|x|\)只能在不跨越零点的区间使用Newton--Leibniz公式，不能跨奇点直接相减。


% ITEM: calculus-example-085
\solutionitem{calculus-example-085}{例5.85 *}
讨论\(\int_0^{+\infty}dx/[x^\alpha(1+x)]\)的敛散性。分别在零点附近与无穷远处比较幂函数；当\(\alpha=1/2\)时，用\(x=t^2\)求积分值。再讨论\(\int_0^1dx/[x^\alpha(1-x)^\beta]\)的敛散条件。
\solutionanswer
在\((0,1]\)，被积函数与\(x^{-\alpha}\)的比介于\(1/2,1\)，零点端要求\(\alpha<1\)；在\([1,+\infty)\)，它与\(x^{-\alpha-1}\)的比介于\(1/2,1\)，无穷端要求\(\alpha>0\)。故恰在\(0<\alpha<1\)收敛。\(\alpha=1/2\)时令\(x=t^2\)，积分\(2\int_0^\infty dt/(1+t^2)=\pi\)。最后一式在两端分别等价于\(x^{-\alpha}\)、\((1-x)^{-\beta}\)，敛散条件为\(\alpha<1,\beta<1\)。


% ITEM: calculus-exercise-59-01
\solutionitem{calculus-exercise-59-01}{习题5.59（1） *}
判断\(\int_1^{+\infty}dx/(x^2+\sqrt{x})\)及\(\int_0^1dx/\sqrt{x+x^2}\)的敛散性，指出比较对象。
\solutionanswer
第一式在\([1,+\infty)\)介于\(0,1/x^2\)，故收敛；第二式在\((0,1]\)介于\(0,1/\sqrt x\)，故收敛。比较函数分别是无穷端的二次幂和零点端的半次幂。


% ITEM: calculus-exercise-59-02
\solutionitem{calculus-exercise-59-02}{习题5.59（2） *}
证明\(\int_0^{+\infty}e^{-x}|\cos x|\,dx\)收敛，进而说明\(\int_0^{+\infty}e^{-x}\cos x\,dx\)收敛。
\solutionanswer
\(0\le e^{-x}|\cos x|\le e^{-x}\)，而后者半轴积分为\(1\)，故绝对值积分收敛，从而原积分绝对收敛。


% ITEM: calculus-exercise-59-03
\solutionitem{calculus-exercise-59-03}{习题5.59（3） *}
将上述分部积分用于\(\int_1^{+\infty}\cos x/x\,dx\)，判断它是否收敛、是否绝对收敛。
\solutionanswer
分部得\(\int_1^R\cos x/x\,dx=[\sin x/x]_1^R+\int_1^R\sin x/x^2dx\)。边界项有极限，后一个积分绝对收敛，故原式收敛。在每个充分靠后的周期中选取\(|\cos x|\ge1/2\)的固定长度子区间，每段绝对值积分下界为正常数乘\(1/n\)，总和发散。因此是条件收敛。


% ITEM: calculus-exercise-60-01
\solutionitem{calculus-exercise-60-01}{习题5.60（1） *}
估计\(\sum_{n=N+1}^{\infty}1/n^2\)。要使尾和小于\(10^{-3}\)，可以选取怎样的\(N\)？
\solutionanswer
单调积分比较给出\(1/(N+1)<\sum_{n=N+1}^\infty n^{-2}<1/N\)。取\(N=1000\)即保证尾和严格小于\(10^{-3}\)。


% ITEM: calculus-exercise-60-02
\solutionitem{calculus-exercise-60-02}{习题5.60（2） *}
讨论\(\sum_{n=2}^{\infty}1/[n(\ln n)^p]\)的敛散性，并说明被积函数在哪一段起单调不增。
\solutionanswer
令\(f(x)=1/[x(\ln x)^p]\)，导数符号由\(-\ln x-p\)决定，所以从\(x\ge\max\{2,e^{-p}\}\)起不增。令\(u=\ln x\)，积分化为\(\int u^{-p}du\)，故级数在\(p>1\)收敛，在\(p\le1\)发散。有限段的增减不影响判别。


% ITEM: calculus-exercise-60-03
\solutionitem{calculus-exercise-60-03}{习题5.60（3） *}
给出\(\int_1^{+\infty}\sin x/x\,dx\)截断到\(R\)的误差界。与绝对收敛情形的尾部比较方法有何不同？
\solutionanswer
分部得\(\int_R^\infty\sin x/x\,dx=\cos R/R-\int_R^\infty\cos x/x^2dx\)，故截断误差不超过\(2/R\)。这里利用振荡的抵消，分部后再控制绝对收敛的余项；原被积函数的绝对值尾部积分本身发散，不能直接拿来作有限误差界。


% ITEM: calculus-exercise-60-04
\solutionitem{calculus-exercise-60-04}{习题5.60（4） *}
若要用数值积分计算\(\int_0^{+\infty}e^{-x^2}\,dx\)，设计一种将总误差分成有限区间求积误差与尾部误差的方案。
\solutionanswer
给定\(\varepsilon>0\)，先选\(R\ge1\)使\(e^{-R^2}/(2R)<\varepsilon/2\)。因为\(x\ge R\)时\(e^{-x^2}\le(x/R)e^{-x^2}\)，尾部满足此界。再在\([0,R]\)用中点公式，取\(M=4R^2+2\)作为二阶导数绝对值的上界，选\(n^2>MR^3/(12\varepsilon)\)，使求积误差小于\(\varepsilon/2\)。两项相加即控制总误差。


% ITEM: calculus-exercise-60-05
\solutionitem{calculus-exercise-60-05}{习题5.60（5） *}
计算\(\int_e^{+\infty}dx/[x(\ln x)^2]\)，并给出级数\(\sum_{n=2}^{\infty}1/[n(\ln n)^2]\)在第\(N\)项之后的尾和上下界。要保证尾和小于\(1/10\)，可以取怎样的\(N\)？比较它与\(\sum1/n^2\)的截断速度。
\solutionanswer
令\(u=\ln x\)，积分为\(\int_1^\infty u^{-2}du=1\)。尾和满足\(1/\ln(N+1)<\sum_{n=N+1}^\infty1/[n(\ln n)^2]<1/\ln N\)。取整数\(N\ge22027>e^{10}\)足以使尾和小于\(1/10\)。它的尾部按\(1/\ln N\)衰减，远慢于\(\sum1/n^2\)的\(1/N\)。


% ITEM: calculus-exercise-60-06
\solutionitem{calculus-exercise-60-06}{习题5.60（6） *}
讨论\(\int_0^1x^{-p}\ln x\,dx\)的敛散性，收敛时求其值。若从\(\varepsilon\)开始积分，求遗漏\((0,\varepsilon)\)所造成的误差，并说明它怎样随\(\varepsilon\to0^+\)趋零。
\solutionanswer
令\(q=1-p\)。\(p<1\)时\(q>0\)，分部或令\(x=e^{-t}\)得到积分\(-1/q^2\)。遗漏区间的积分绝对值为\(\varepsilon^q[-\ln\varepsilon/q+1/q^2]\to0\)。\(p\ge1\)时，令\(x=e^{-t}\)可见负的积分发散至\(-\infty\)，故不收敛。


% ITEM: calculus-exercise-60-07
\solutionitem{calculus-exercise-60-07}{习题5.60（7） *}
设\(p>0\)。用分部积分证明\(\int_1^{+\infty}\sin x/x^p\,dx\)收敛，并给出截断到\(R\)的误差界。再按\(p\)的范围判断绝对收敛与条件收敛；讨论绝对值时，可以在每个周期中选取\(|\sin x|\ge1/2\)的子区间。
\solutionanswer
分部有\(\int_R^\infty\sin x/x^pdx=\cos R/R^p-p\int_R^\infty\cos x/x^{p+1}dx\)，故对每个\(p>0\)收敛，尾部绝对值不超过\(2/R^p\)。\(p>1\)时绝对收敛；\(0<p\le1\)时，每周期取\(|\sin x|\ge1/2\)的固定长度子区间，绝对值积分下界为\(C/n^p\)，累加发散，因此条件收敛。


% ITEM: calculus-exercise-60-08
\solutionitem{calculus-exercise-60-08}{习题5.60（8） *}
正文的对比例子为\(g_n(x)=n^{-1}(1-x/n)\)（\(0\le x\le n\)）、其余为零：在整个半轴一致趋零，而半轴积分均为\(1/2\)。
设\(f_n(x)=n^2xe^{-nx}\)，\(x\ge0\)。求逐点极限及每个\(f_n\)在半轴上的积分。说明为何不能直接交换积分与极限，并比较它与正文中图像向远处铺开的例子：本题的面积集中到了哪里？
\solutionanswer
每个固定\(x>0\)有\(n^2xe^{-nx}\to0\)，零点也为零，逐点极限为零。令\(u=nx\)，每个半轴积分却均为\(\int_0^\infty ue^{-u}du=1\)。故不能交换极限和积分。最大值在\(x=1/n\)，高度为\(n/e\)，面积集中到越来越靠近零点、越来越窄的区域，而非移向无穷远。
\par
